%Dans une station balnéaire on a interrogé 600 touristes, français ou étrangers, sur leur séjour.
%Tous ont répondu. Les effectifs sont reportés dans le tableau suivant :
%
%\begin{tblr}	{
%    colspec = {*5{X[c,m]}},
%	  hline{2-Z},
%    vline{2-Z},
%    vline{1} = {2-4}{},
%    hline{1} = {2-5}{},
%	}
%	& Camping 	& Location 	& Hôtel 	& \textbf{Total} 	\\
%Français 			& 232 		& 116 		& 12 	& 360 	\\
%Étrangers 			& 96 		& 96 		& 48 	& 240 	\\
%\textbf{Total} 		& 328 		& 212 		& 60 	& 600 	\\
%\end{tblr}
%
%Les réponses des questions suivantes seront données sous forme de fractions non nécessairement simplifiées.
%
%On choisit une personne au hasard. On suppose que toutes les personnes ont la même probabilité d’être choisies.
%
%On considère les événements :
%\begin{itemize}
%\item $E$ : « La personne choisie est un touriste étranger ».
%\item $C$ : « La personne choisie séjourne dans un camping ».
%\end{itemize}
%
Puisque la personne est choisie au hasard, on est dans une situation d'équiprobabilité.
\begin{enumerate}
  \item 
  %Donner la probabilité $P\left({C}\right)$ de l’événement $C$.
  Il y a 600 touristes, dont 328 séjournent dans un camping, donc \mbox{$P(C)=\dfrac{328}{600}=\dfrac{41}{75}$}
  \item 
  %Donner la probabilité $P\left(\overline{C}\right)$ où $\overline{C}$ désigne l’événement contraire de $C$.
  $\overline{C}$ est l’événement contraire de $C$, donc \mbox{$P(\overline{C})=1-P(C)=\dfrac{600}{600}-\dfrac{328}{600}=\dfrac{272}{600}=\dfrac{34}{75}$}.
  \item 
  %Donner la probabilité de l'événement : « La personne interrogée est un touriste étranger et séjourne dans un camping ».
  On cherche \mbox{$P(E\cap C)$}. Il y a 600 personnes, dont 96 sont à la fois des touristes étrangers et séjournent dans un camping.
  
  Donc \mbox{$P(E\cap C)=\dfrac{96}{600}=\dfrac{4}{25}$}.
  \item
%  Sachant que la personne choisie est en location, donner la probabilité qu'elle soit un touriste français.
   Il y a 212 personnes qui sont en location, dont 116 français. Donc la probabilité que la personne choisie, sachant qu'elle est en location, soit française, est \mbox{$\dfrac{116}{212}=\dfrac{29}{53}$}.
\end{enumerate}


\needspace{10\baselineskip}
% saut de page si moins de 10 lignes disponibles
